\section{Dependencies on prior sequence papers}
\label{app:dependencies}

The paper relies on results from nine prior sequence papers
(\citep{lasser2026gfm,lasser2026poset,lasser2026horizon,lasser2026scm,lasser2026exo,lasser2026phase,lasser2026wamura,lasser2026paper8a,lasser2026paper8b}).
Table~\ref{tab:dependencies} summarises which result is used
where.

\begin{table}[ht]
\centering
\footnotesize
\setlength{\tabcolsep}{4pt}
\begin{tabular}{>{\raggedright\arraybackslash}p{0.18\linewidth}
                >{\raggedright\arraybackslash}p{0.32\linewidth}
                >{\raggedright\arraybackslash}p{0.42\linewidth}}
\toprule
\textbf{Source paper} & \textbf{Result imported} &
\textbf{Role in this paper} \\
\midrule
\citep[Foundation]{lasser2026gfm} &
$\volL$ as poset measure (Definition~6); naming of B-to-C gap
(Section~7) &
The proxy whose proxy-truth gap the paper bounds. \\
\addlinespace
\citep[Poset Definition]{lasser2026poset} &
Axioms M1--M6 (Proposition~1); $\volL$ self-balancing
(Theorem~1) &
The structural apparatus utility theory lacks; M5/M6 carry the
non-fungibility, cooperative, individuation features beyond
$\Pstd$. \\
\addlinespace
\citep[Horizon Aware]{lasser2026horizon} &
Anti-monopolar property (Proposition~6); $\gamma^*$ threshold;
exercise pathway (Definition~2); preemptive restriction
(Proposition~1) &
Structural prediction with no utility-theoretic analogue;
risk machinery the Goodhart theorem perturbs against. \\
\addlinespace
\citep[SCM]{lasser2026scm} &
Causal-structure formalism &
Substrate for the optimization-pressure argument in
Conjecture~\ref{conj:optimization_pressure}. \\
\addlinespace
\citep[Exogenous Verification]{lasser2026exo} &
Wireheading discussion; commitment protocol &
Goodhart preview (the original wireheading note); commitment
layer composes with~\citep{lasser2026paper8a}'s sacrifice
channel. \\
\addlinespace
\citep[Phase Redundancy]{lasser2026phase} &
Phase boundary (Theorem~1); design criterion (Theorem~2) &
Alignment property whose perturbation under variable proxy is
Open Question~\ref{oq:phase_boundary}. \\
\addlinespace
\citep[Controlled Relaxation]{lasser2026wamura} &
Convergence (Theorem~2); damage bound (Theorem~1) &
Alignment properties listed in the perturbative roadmap
(Section~\ref{sec:goodhart_examples}). \\
\addlinespace
\citep[Revealed Sacrifice]{lasser2026paper8a} &
Aggregate B-to-C lower bound (Theorem~2); monotone accumulation
(Propositions~2~and~3); event-local bundle decomposition
(Definition~4); residual-class structure; assumptions
S0--S5; Remark~18 (external attestation) &
The measurement channel that drives $\betaL$ upward; the formal
S0 bridge from inner preference to operational truth; the
$\alpha_{n,c}$ substrate for Channel~4; the attestation
discipline for Channel~2. \\
\addlinespace
\citep[Need Sufficiency]{lasser2026paper8b} &
Need-sufficiency architecture (Definition~2); polarity boundary
(Definition~3); gap decomposition (Definition~8);
wireheading-consistent HHI (Propositions~6~and~7); residual class
$\ResS$ (Definition~12); need-access cost (Definition~9) &
The diagnostic architecture that interprets $\betaL$; the
$\delta$-decomposition cells through which the four channels
of Section~\ref{sec:lineage} act; the polarity boundary that
distinguishes need-share from want-share trades; $\ResS$ as
the formal version of the welfare-economics shadow-economy
gap. \\
\bottomrule
\end{tabular}
\caption{Dependencies on prior sequence papers and their roles
in this paper.  All references are pinned to current preprint
formulations; if the upstream preprints revise their numbering
or content, downstream references will need refreshing
correspondingly.}
\label{tab:dependencies}
\end{table}

The capability-welfare-economics references
(Sen~1985~\citep{sen1985commodities},
Nussbaum~2000~\citep{nussbaum2000women},
Becker~1965~\citep{becker1965theory},
Stiglitz--Sen--Fitoussi~2009~\citep{stiglitz2009report},
Arrow~1951~\citep{arrow1951social}) are positioning rather
than load-bearing in the formal sense: their qualitative
content is what the GFM apparatus formalizes, and the
formalization-claim does not technically depend on any specific
result from these references.  Section~\ref{sec:lineage_historical}
makes the structural correspondences explicit.
